\chapter{Introduction}

The dynamics of firm growth and economic fluctuations have long been modeled through the lens of complex interacting systems. One of the most robust empirical observations in quantitative economics—often referred to as a "stylized fact"—is the scaling relationship between the size of a firm ($S$) and the standard deviation of its growth rate ($\sigma$). Empirical data consistently demonstrates a power-law decay of the form $\sigma(S) \propto S^{-\beta}$, where the scaling exponent $\beta$ is typically found in the range of $0.15$ to $0.20$. From a statistical physics perspective, this exponent is highly anomalous. If a firm were composed of independent, uncorrelated sub-units, the Central Limit Theorem (CLT) would strictly dictate an exponent of $\beta = 0.5$. The empirically observed suppression of this exponent implies the existence of deep, systemic, intra- and inter-firm correlations that cannot be explained by standard additive random walks.

To uncover the microscopic origins of these macroscopic correlations, agent-based network models have become a primary theoretical tool. Specifically, the Generalized Lotka-Volterra (GLV) model—adapted from theoretical ecology—provides a deterministic framework to study the competitive and cooperative interactions between distinct economic nodes. However, extracting continuous, non-decaying volatility from a deterministic model requires situating the system at a very specific thermodynamic boundary: the critical phase transition ($\mu = \mu_c$). At topological criticality, the restoring forces of the system vanish, and long-lived transient macroscopic fluctuations naturally emerge.

Simulating interacting systems exactly at this divergence point introduces a severe computational bottleneck. As the interaction strength approaches the critical threshold, the macroscopic variables diverge, causing standard continuous ordinary differential equation (ODE) solvers to fail due to infinitesimally small step sizes or machine precision overflow.

The primary objective of this thesis is twofold. First, we engineer a robust numerical framework to bypass this critical divergence. By applying a coordinate transformation that maps the unbounded abundances onto a bounded probability simplex, coupled with a dynamic time-rescaling, we establish a method to integrate the GLV equations seamlessly through the critical singularity. Second, we deploy this computational tool to formally test whether the transient critical dynamics of the static GLV model are sufficient to break the Central Limit Theorem and recover the empirical volatility scaling exponent of $\beta \approx 0.2$.

Applying this stabilized numerical framework, we conduct a rigorous stress-test of the GLV model across different network topologies to evaluate its capacity to generate realistic economic fluctuations. We first investigate the system on networks with exponential degree distributions. Because the variance of an exponential degree distribution is finite and well-behaved, a macroscopic critical threshold ($\mu_c > 0$) exists, allowing the system to be simulated at criticality. However, our findings indicate that the resulting volatility scaling exponent is tightly bound around $\beta \approx 0.5$. This mathematically demonstrates that, in the absence of heavy-tailed topological hierarchies or multiplicative external shocks, the deterministic interactions within the static GLV model effectively average out as mean-field noise. Consequently, the system remains strictly governed by the Central Limit Theorem, failing to produce the deep correlations required to match empirical economic data.

Furthermore, we extend our analysis to scale-free (power-law) networks, which more accurately reflect the topology of real-world economic and financial systems. In this regime, the theoretical limits of the continuous GLV model are laid bare. Because the second moment of a power-law degree distribution diverges in the thermodynamic limit ($\langle g^2 \rangle \to \infty$), the theoretical critical threshold vanishes ($\mu_c \to 0$). We demonstrate that attempting to simulate such a system does not merely result in a failure of the algorithm, but reflects a physical reality of the model: an instantaneous "monopoly collapse." A single massive hub node exponentially dominates the system, forcing the relative fractions of all other nodes in the simplex to drop so violently that they trigger a 64-bit machine precision underflow. This confirms that without structural cutoffs or artificial bounds, continuous GLV interactions are incompatible with the scale-free topologies required to break the Central Limit Theorem.

Ultimately, this thesis maps the exact boundaries of where static topological criticality ceases to describe economic reality, providing a definitive falsification of the simple GLV model as a standalone mechanism for empirical volatility scaling.

The remainder of this thesis is structured as follows. Chapter 2 details the mathematical derivation of the coordinate transformation, formally mapping the diverging GLV equations to the probability simplex, and discusses the limits of the analytical mean-field approximations. Chapter 3 presents the numerical methodology and the empirical results, illustrating the successful bypass of the computational divergence and analyzing the resulting scaling behaviors on both exponential and power-law networks. Finally, Chapter 4 concludes the work, summarizing the implications of these findings and outlining potential theoretical extensions, such as the introduction of multiplicative stochastic noise, required to recover realistic macroeconomic dynamics.
